% -*- coding: utf-8 -*-
% !TEX program = xelatex
\documentclass[plain,noanswer,medmath]{../../template/jnuexam}
\usepackage{../../template/kysx}

\begin{document}


\examtitle{1987}{1}


\makepart{填空题}{本题满分15分，每小题3分}


\begin{problem}
与两直线$\begin{cases}
  x = 1 \\
  y = - 1 + t \\
  z = 2 + t
\end{cases}$
及$\frac{x + 1}{1} = \frac{y + 2}{2} = \frac{z - 1}{1}$都平行，且过原点的平面方程为 \fillin{}．
\end{problem}


\begin{problem}
当$x =$      时，函数$y = x2^x$取得极小值 \fillin{}．
\end{problem}


\begin{problem}
由曲线$y = \ln x$与两直线$y = (\e + 1) - x$及$y = 0$所围成的平面图形的面积是 \fillin{}．
\end{problem}


\begin{problem}
设$L$为取正向的圆周$x^2 + y^2 = 9$，则曲线积分$\oint_L (2xy - 2y)\dx + (x^2 - 4x)\dy$的值是 \fillin{}．
\end{problem}


\begin{problem}
已知三维线性空间的一组基底为$\alpha_1 = (1,1,0)^T$,$\alpha_2 = (1,0,1)^T$,$\alpha_3 = (0,1,1)^T$,
则向量$\beta = (2,0,0)^T$在上述基底下的坐标是 \fillin{}．
\end{problem}


\makepart{}{本题满分8分}

\begin{problem}
求正的常数$a$与$b$，使等式
$\smash[t]{\lim_{x \to 0}\frac{1}{bx-\sin x}\int_0^x \frac{t^2}{\sqrt{a + t^2}}\dt = 1}$成立．
\end{problem}


\makepart{计算题}{本题满分7分}


\begin{problem}[points=3]
设$f$, $g$为连续可微函数，$u = f(x,xy)$, $v = g(x+xy)$,
求$\frac{\pdu}{\pdx}$，$\frac{\pdv}{\pdx}$．
\end{problem}


\begin{problem}[points=4]
设矩阵$A$和$B$满足关系式$AB = A + 2B$，其中$\smash[t]{A=\begin{pmatrix}
 3&0&1 \\
 1&1&0 \\
 0&1&4
\end{pmatrix}}$，求矩阵$B$．
\end{problem}


\makepart{}{本题满分8分}

\begin{problem}
求微分方程$y''' + 6y'' + (9 + a^2)y' = 1$的通解（一般解），其中常数$a > 0$．
\end{problem}


\makepart{选择题}{本题满分12分，每小题3分}


\begin{problem}
设常数$k > 0$，则级数$\sum_{n = 1}^{\infty}(- 1)^n\frac{k + n}{n^2}$\pickout{}
\begin{abcd}
\item 发散．						
\item 绝对收敛．
\item 条件收敛．				
\item 收敛或发散与k的取值有关．
\end{abcd}
\end{problem}


\begin{problem}
设$f(x)$为已知连续函数，$I=t\int_0^{s/t}\!\!f(tx)\dx$，其中$s>0$, $t>0$，则$I$的值\pickout{}
\begin{abcd}
\item 依赖于$s$和$t$．
\item 依赖于$s$,$t$,$x$．
\item 依赖于$t$和$x$，不依赖于$s$．
\item 依赖于$s$，不依赖于$t$．
\end{abcd}
\end{problem}


\begin{problem}
设$\lim_{x \to a} \frac{f(x)-f(a)}{(x-a)^2}=-1$，则在点$x=a$处\pickout{}
\begin{abcd}
\item $f(x)$的导数存在，且$f'(a) \ne a$．			
\item $f(x)$取得极大值．
\item $f(x)$取得极小值．					
\item $f(x)$的导数不存在．
\end{abcd}
\end{problem}


\begin{problem}
设$A$为$n$阶方阵，且$A$的行列式$|A|=a \ne 0$，而$A^*$为$A$的伴随矩阵，则$|A^*|$等于\pickout{}
\begin{abcd}
\item $a$．       
\item $\frac{1}{a}$．     
\item $a^{n - 1}$．
\item $a$．
\end{abcd}
\end{problem}


\makepart{}{本题满分10分}

\begin{problem}
求幂级数$\sum_{n = 1}^{\infty}\frac{1}{n2^n}x^{n + 1}$的收敛域，并求其和函数．
\end{problem}


\makepart{}{本题满分10分}

\begin{problem}
计算曲面积分
$$I = \iint_{\Sigma}x(8y + 1)\dy\dz + 2(1 - y^2)\dz\dx - 4yz\dx\dy,$$
其中$\Sigma$是由曲线$\begin{cases}
  z = \sqrt{y - 1} \\
  x = 0
\end{cases}$（$1 \le y \le 3$）
绕$y$轴旋转一周所形成的曲面，它的法向量与$y$轴正向的夹角恒大于$\frac{\pi}{2}.$
\end{problem}


\makepart{}{本题满分10分}

\begin{problem}
设函数$f(x)$在闭区间$[0,1]$上可微，对于$[0,1]$上的每一个$x$，
函数$f(x)$的值都在开区间$(0,1)$内，且$f'(x) \ne 1$，
证明在$(0,1)$内有且仅有一个$x$，使$f(x) = x$．
\end{problem}


\makepart{}{本题满分8分}

\begin{problem}
问$a$, $b$为何值时，线性方程组$\begin{cases}
  x_1 + x_2 + x_3 + x_4 = 0, \\
  x_2 + 2x_3 + 2x_4 = 1, \\
  -x_2 + (a - 3)x_3 - 2x_4 = b, \\
  3x_1 + 2x_2 + x_3 + ax_4 = - 1
\end{cases}$有唯一解、无解、有无穷多组解？并求出有无穷多组解时的通解．
\end{problem}


\makepart{填空题}{本题满分6分，每小题2分}


\begin{problem}
设在一次试验中事件$A$发生的概率为$p$，现进行$n$次独立试验，
则$A$至少发生一次的概率为 \fillin{}；而事件$A$至多发生一次的概率为 \fillin{}．
\end{problem}


\begin{problem}
三个箱子，第一个箱子中有$4$个黑球$1$个白球，第二个箱子中有$3$个黑球$3$个白球，
第三个箱子中有$3$个黑球$5$个白球．现随机地取一个箱子，再从这个箱子中取出$1$个球，
这个球为白球的概率等于 \fillin{}．已知取出的球是白球，此球属于第二个箱子的概率为 \fillin{}．
\end{problem}


\begin{problem}
已知连续随机变量$X$的概率密度为$f(x)=\frac{1}{\sqrt{\pi}}\e^{-x^2+2x-1}$，
则$X$的数学期望为 \fillin{}；$X$的方差为 \fillin{}．
\end{problem}


\makepart{}{本题满分6分}

\begin{problem}
设随机变量$X$, $Y$相互独立，其概率密度函数分别为
$$f_X(x) = \begin{cases}
  1, & 0 \le x \le 1, \\
  0, & \text{其它},
\end{cases}\qquad
f_Y(y) = \begin{cases}
  \e^{-y}, & y > 0, \\
        0, & y \le 0,
\end{cases}$$
求随机变量$Z = 2X + Y$的概率密度函数．
\end{problem}


\end{document}
